\documentclass[11pt,a4paper]{article}
\usepackage[utf8]{inputenc}
\usepackage[T1]{fontenc}
\usepackage{geometry}
\geometry{margin=2.5cm}
\usepackage{booktabs}
\usepackage{hyperref}
\usepackage{enumitem}

\title{Music History \& Heritage Knowledge Graph\\[0.3em]\large Knowledge Engineering CW2 Report}
\author{Lucas Perez Reis Lobo, Yutaka Shibata, Sofia Davis, \\ William Caulton, Mert Yerlikaya, Davyd Shtepa}
\date{April 2026}

\begin{document}
\maketitle

\section{Introduction}

This project constructs an automated knowledge graph for the Music History domain, covering 58 primary artists from medieval (Hildegard von Bingen, 1098) to modern (BTS, Kendrick Lamar). The system extends two existing ontologies (Music Ontology and Schema.org) and populates them with data from 4 sources: MusicBrainz, Discogs, Wikidata (structured), and Wikipedia (textual). The entire process is automated through a Python pipeline that fetches data from APIs, maps it to RDF, performs entity linking, validates triples, ingests RAG completion results, and asserts defined class instances. The final KG contains 42,554 triples covering 1,332 artists, 1,408 releases, 860 musical works, 274 genres, 257 awards, and 62 instruments. All API responses and LLM extraction results are cached, enabling reproducible builds in under 35 seconds. Team: Sofia Davis and Davyd Shtepa (domain experts), William Caulton and Mert Yerlikaya (modelling experts), Lucas Perez Reis Lobo and Yutaka Shibata (requirements and pipeline).

\section{Data Source Selection}

Four sources were selected from 6 evaluated across 5 test artists spanning rock, jazz, classical, Brazilian, and African music. MusicBrainz serves as the primary hub, providing artist metadata, releases, tracks, and artist-to-artist relationships (membership, influences, collaborations, production credits). Every entity has a UUID usable as a URI, and URL relationships enable cross-referencing to all other sources. Discogs supplements with a genre-to-style hierarchy that maps directly to our \texttt{subgenreOf} property. Wikidata provides awards, instruments, influences, and record labels as structured RDF, eliminating the need for LLM extraction of these properties. Its class hierarchy (P31/P279*) is used to programmatically filter music genres and music awards. Wikipedia provides narrative text for LLM-based triple extraction, capturing facts unavailable in structured sources: alter egos, album groupings, narrative genres, and producer relationships. DBpedia was rejected (redundant with Wikidata, noisy template parameters), Open Opus was rejected (classical-only), and Last.fm was inaccessible. For 6 artists where MusicBrainz name-based search returned incorrect results (including Umm Kulthum due to Arabic script), MBID overrides are configured.

\section{Extension of Existing Ontologies}

Music Ontology was extended with 2 subclasses (\texttt{CoverRecording} subClassOf \texttt{mo:Signal} for the 373 cover recordings identified by cross-referencing MusicBrainz work-rels with Wikidata composers; \texttt{MusicalWork} subClassOf \texttt{mo:MusicalWork} distinguishing abstract compositions from recordings) and 2 subproperties (\texttt{collaboratedWith} subPropertyOf \texttt{mo:collaborated\_with}, declared symmetric; \texttt{playsInstrument} subPropertyOf \texttt{mo:instrument}). Schema.org was extended with 2 subclasses (\texttt{MultinationalBand} subClassOf \texttt{schema:MusicGroup} with 5 instances asserted via SPARQL; \texttt{MusicAward} subClassOf \texttt{schema:Award} filtering non-music awards) and 2 subproperties (\texttt{signedTo} subPropertyOf \texttt{schema:affiliation}; \texttt{countryOfOrigin} subPropertyOf \texttt{schema:nationality}). External classes and properties are declared as minimal stubs without domain/range to avoid conflicting with the original ontology definitions.

Three defined classes use \texttt{owl:equivalentClass}: \texttt{AwardWinningArtist} (MusicArtist AND wonAward some MusicAward, 52 instances), \texttt{CollaboratingArtist} (MusicArtist AND collaboratedWith some MusicArtist, 118 instances), and \texttt{ProducerArtist} (MusicArtist AND produced some Release, 10 instances). \texttt{InternationalCollaborator} is a subclass of \texttt{CollaboratingArtist} with 11 instances asserted programmatically, since OWL2 cannot compare property values across individuals. Under the Open World Assumption, absence of a \texttt{wonAward} triple does not entail an artist has not won awards; the pipeline therefore asserts defined class membership explicitly via SPARQL. The ontology declares 21 OWL2 features: 1 symmetric, 1 transitive, 7 irreflexive, 4 asymmetric, 6 inverse pairs (all materialised), and 2 disjointness axioms.

\section{Mappings}

The pipeline automates 15 steps: ontology header generation, LLM text extraction via Google Gemini API, per-artist data fetching from 4 APIs, structured mapping (JSON to RDF via rdflib), text mapping (LLM extraction JSON to RDF with 5-tier entity linking at 99\% accuracy), enrichment of secondary artists with genres and dates, URI consolidation merging duplicate entities, type assignment for orphan entities, cover recording detection, multinational band and international collaborator classification, validation via 22 structural rules, RAG ingestion of 494 completion triples, and defined class assertion. Course techniques demonstrated include R2RML mapping documents (Week~8), SPARQL Anything CONSTRUCT queries (Week~8), and LLM-based triple extraction (Week~7). Text extraction uses a controlled predicate vocabulary with a ``Do NOT extract'' list preventing duplication with Wikidata. The template was applied to all 58 artists (897 triples, 44\% using text-unique predicates). Entity linking uses a typed label index with dual-scorer gate, stopword removal, and subject type inference from predicates. All 22 validation rules are pattern-based with zero hardcoded entity-specific fixes, addressing issues from self-references and type mismatches to disjointness violations and orphan genres.

\section{Queries}

Twenty competency questions guide the ontology design. The 10 manual CQs were developed by the domain and modelling experts: CQ1 album release dates, CQ2 artist genres, CQ3 track-to-album lookup, CQ4 album producers, CQ5 band members, CQ6 producers by genre, CQ7 instruments in bands, CQ8 award counts, CQ9 subgenre hierarchy, CQ10 composition dates with recording gap. The 10 LLM-augmented CQs fill gaps identified in the manual set: CQ11 award-winning artists at labels, CQ12 influences sharing genres, CQ13 label genre diversity, CQ14 cross-country collaborations, CQ15 instruments per genre, CQ16 multinational bands, CQ17 artist label discography, CQ18 cross-country cover recordings (4-entity join), CQ19 collaborating award winners, CQ20 genre geographic spread by decade.

The LLM CQs were generated through iterative prompting (P3--P6) using role-based, few-shot, and constrained output techniques. All 20 are answered by SPARQL queries using patterns including \texttt{GROUP\_CONCAT}, \texttt{COUNT} with \texttt{HAVING}, property paths for transitive closure (\texttt{subgenreOf+}), 4-entity joins, and self-joins with deduplication filters. Counter-examples demonstrate correct exclusion: Hermeto Pascoal has 0 awards (CQ8), The Beatles are single-country (CQ16), same-country collaborators are filtered (CQ14), and bebop is a leaf genre (CQ9). Known limitations include CQ17's cartesian product from unscoped \texttt{signedTo} and CQ10's time gap only computable for works with cover recordings.

\section{Evaluation Methodology}

Performance: the pipeline completes in 33.9 seconds from cache, with 47.6~MB peak memory. Quality was measured using CM1--CM4 completeness metrics \cite{zaveri2015}: CM1 schema completeness reached 97\% (29/30 properties populated), CM2 property completeness for primary artists is high (country 98\%, genre 100\%), and all 20 CQs return results (100\% coverage). Three KG embedding models were trained with PyKEEN: RotatE achieved MRR 0.557, CompGCN 0.526, TransE 0.086. Materialising inverse triples improved MRR from 0.05 to 0.55 (11$\times$ improvement), validating the ontology decision to materialise all 6 inverse pairs. Data quality cleanup shifted the optimal model from CompGCN to RotatE, demonstrating that graph sparsity determines which architecture performs best. NER comparison (Week~7) showed SpaCy finds the most entities but only LLM extraction produces structured relations. Six RAG vs Plain LLM comparisons demonstrated that RAG functions as a data quality auditor, surfacing non-music genres, instrument duplicates, and range violations that led to 6 pipeline fixes. The iterative cycle of build, evaluate with RAG, fix, and rebuild proved more effective than either pure automation or manual curation. A domain expert review identified 76 non-music genres and 31 non-music awards, all fixed structurally. The KG was validated in Prot\'{e}g\'{e} for ontological correctness.

\section{Repository Link}

\url{https://github.com/LucasPRLobo/ke-cw2}

Full technical details including pipeline code, experiment notebooks, evaluation scripts, cached data, and the complete prompts log are available in the repository. While the pipeline fully automates KG construction, three categories of manual intervention remain essential: ontological correctness (verified in Prot\'{e}g\'{e}), semantic accuracy (domain expert review), and multilingual entity resolution (MBID overrides for non-Latin scripts).

\begin{thebibliography}{1}
\bibitem{zaveri2015} Zaveri, A., Rula, A., Maurino, A., Pietrobon, R., Lehmann, J., \& Auer, S. (2016). Quality assessment for Linked Data: A survey. \textit{Semantic Web}, 7(1), 63--93.
\end{thebibliography}

\end{document}
